\section{CONCLUSION}
\label{sec:conclusion}

In this work, we presented a comprehensive investigation into RT-FedMTL using SLMs across diverse NLU benchmarks. Our systematic evaluation across five architectures provided critical insights for the deployment of decentralized NLU systems. Our results demonstrate that FL effectively preserves data privacy while significantly optimizing resource distribution, achieving up to 93\% reduction in per-client GPU resource consumption. We identified a critical performance-efficiency trade-off: DistilBERT consistently offered the most robust performance for sentiment analysis, showing superior resilience to negative transfer in SST-2. However, for a unified multi-task suite, MiniLM emerged as the most versatile and efficient architecture, outperforming larger models in QQP and STS-B tasks despite being approximately 3$\times$ smaller in memory footprint (0.38\,GB vs 1.09\,GB). For scenarios where memory and bandwidth are the primary bottlenecks, architectures like MiniLM provide the optimal balance of stability and throughput for real-time federated multi-task environments.

Future research will focus on developing adaptive task-sharing strategies to further mitigate negative transfer in multi-head architectures. Additionally, we plan to explore personalized federated learning techniques and the integration of differentially private mechanisms within the real-time communication layer, ensuring that efficiency gains do not compromise privacy in production NLU deployments.
